% ============================================================
%  Programação em Java — Lição 3: Listas, Arrays e Strings
%  Java Programming MOOC · Universidade de Helsinque
%  Compilar com: pdflatex licao03.tex
% ============================================================
\documentclass[12pt, a4paper]{article}

% ── Codificação e língua ──────────────────────────────────
\usepackage[utf8]{inputenc}
\usepackage[T1]{fontenc}
\usepackage[brazil]{babel}

% ── Fontes ────────────────────────────────────────────────
\usepackage{palatino}
\usepackage[scaled=0.88]{beramono}

% ── Geometria ─────────────────────────────────────────────
\usepackage[
  a4paper,
  top=3cm, bottom=3cm,
  left=3cm, right=3cm
]{geometry}

% ── Cores ─────────────────────────────────────────────────
\usepackage[dvipsnames, table, x11names]{xcolor}

\definecolor{javaOrange}{HTML}{E76F00}
\definecolor{javaDeep}{HTML}{BF5700}
\definecolor{javaBlue}{HTML}{1565C0}
\definecolor{javaTeal}{HTML}{00695C}
\definecolor{javaAmber}{HTML}{B45309}
\definecolor{javaInk}{HTML}{2B2140}
\definecolor{javaCodeBg}{HTML}{FFF8F0}
\definecolor{javaRule}{HTML}{FFD7B5}
\definecolor{javaLightBg}{HTML}{FFF5EC}
\definecolor{javaKeyword}{HTML}{D73A49}
\definecolor{javaString}{HTML}{032F62}
\definecolor{javaComment}{HTML}{6A737D}

% ── Cabeçalho e rodapé ────────────────────────────────────
\usepackage{fancyhdr}
\pagestyle{fancy}
\fancyhf{}
\renewcommand{\headrulewidth}{1.5pt}
\renewcommand{\headrule}{\hbox to\headwidth{\color{javaOrange}\leaders\hrule height \headrulewidth\hfill}}
\fancyhead[L]{\small\color{javaOrange}\textbf{Programação em Java}}
\fancyhead[R]{\small\color{javaDeep}Lição 3: Listas, Arrays e Strings}
\fancyfoot[C]{\small\color{javaAmber}\thepage}
\fancyfoot[R]{\small\color{gray}Java Programming MOOC}

% ── Títulos de seção ──────────────────────────────────────
\usepackage{titlesec}

\titleformat{\section}
  {\color{javaOrange}\large\bfseries}
  {\color{javaOrange}\thesection.}{0.6em}{}
  [\vspace{2pt}\color{javaRule}\titlerule]

\titleformat{\subsection}
  {\color{javaDeep}\normalsize\bfseries}
  {\color{javaDeep}\thesubsection.}{0.5em}{}

\titleformat{\subsubsection}
  {\color{javaTeal}\normalsize\bfseries}
  {\color{javaTeal}\thesubsubsection.}{0.5em}{}

\titlespacing*{\section}{0pt}{2ex}{1ex}
\titlespacing*{\subsection}{0pt}{1.8ex}{0.6ex}
\titlespacing*{\subsubsection}{0pt}{1.4ex}{0.4ex}

% ── Hiperlinks ────────────────────────────────────────────
\usepackage[
  colorlinks=true,
  linkcolor=javaBlue,
  urlcolor=javaBlue,
  citecolor=javaBlue,
  pdfauthor={Java Programming MOOC},
  pdftitle={Lição 3: Listas, Arrays e Strings},
  pdfsubject={Programação em Java},
  pdfkeywords={java, listas, arrays, strings, ArrayList}
]{hyperref}

% ── Código-fonte ─────────────────────────────────────────
\usepackage{listings}

\lstdefinestyle{java}{
  language=Java,
  basicstyle=\small\ttfamily\color{javaInk},
  keywordstyle=\color{javaKeyword}\bfseries,
  stringstyle=\color{javaString},
  commentstyle=\color{javaComment}\itshape,
  numberstyle=\tiny\color{gray},
  numbers=left,
  numbersep=8pt,
  stepnumber=1,
  showstringspaces=false,
  breaklines=true,
  breakatwhitespace=true,
  tabsize=4,
  frame=none,
  morekeywords={var, record, String, Scanner, ArrayList, HashMap,
                System, Integer, Double, Boolean, File, Math},
}

\lstdefinestyle{plaintext}{
  basicstyle=\small\ttfamily\color{javaInk},
  numbers=none,
  showstringspaces=false,
  breaklines=true,
  frame=none,
}

% ── tcolorbox para blocos de código ───────────────────────
\usepackage[most]{tcolorbox}
\tcbuselibrary{listings, breakable}

\newtcblisting{javacode}{
  listing only,
  listing style=java,
  breakable,
  enhanced jigsaw,
  colback=javaCodeBg,
  colframe=javaRule,
  borderline west={3pt}{0pt}{javaDeep},
  arc=4pt,
  top=6pt, bottom=6pt, left=10pt, right=6pt,
  boxrule=0.5pt,
}

\newtcblisting{plaincode}{
  listing only,
  listing style=plaintext,
  breakable,
  enhanced jigsaw,
  colback=javaCodeBg,
  colframe=javaRule,
  borderline west={3pt}{0pt}{javaAmber},
  arc=4pt,
  top=6pt, bottom=6pt, left=10pt, right=6pt,
  boxrule=0.5pt,
}

% Caixa de destaque (callout)
\newtcolorbox{callout}[1][Nota]{
  breakable,
  enhanced jigsaw,
  colback=javaLightBg,
  colframe=javaOrange,
  coltitle=white,
  fonttitle=\small\bfseries,
  title=#1,
  arc=4pt,
  boxrule=1pt,
  top=4pt, bottom=4pt, left=8pt, right=8pt,
}

% Caixa de exercício
\newtcolorbox{exercicio}[1]{
  breakable,
  enhanced jigsaw,
  colback=white,
  colframe=javaRule,
  borderline west={4pt}{0pt}{javaOrange},
  coltitle=javaOrange,
  fonttitle=\normalsize\bfseries,
  title=Exercício #1,
  arc=4pt,
  boxrule=0.5pt,
  top=6pt, bottom=6pt, left=10pt, right=8pt,
  attach boxed title to top left={yshift=-2mm, xshift=8pt},
  boxed title style={
    colback=javaOrange, colframe=javaDeep,
    arc=3pt, boxrule=0.5pt,
    fontupper=\small\bfseries\color{white}
  },
}

% ── Tabelas ───────────────────────────────────────────────
\usepackage{booktabs}
\usepackage{longtable}
\usepackage{array}
\usepackage{multirow}

% ── Utilitários ───────────────────────────────────────────
\usepackage{microtype}
\usepackage{setspace}
\usepackage{parskip}
\setlength{\parskip}{0.5em}
\usepackage{enumitem}
\setlist[itemize]{topsep=2pt, itemsep=2pt, leftmargin=1.5em}
\setlist[enumerate]{topsep=2pt, itemsep=2pt, leftmargin=1.5em}

% ── Código inline ─────────────────────────────────────────
\newcommand{\code}[1]{\texttt{\color{javaDeep}#1}}

% ──────────────────────────────────────────────────────────
%  CAPA
% ──────────────────────────────────────────────────────────
\usepackage{eso-pic}
\usepackage{tikz}

\newcommand{\makejavacapa}{%
  \begin{titlepage}
    \begin{tikzpicture}[remember picture, overlay]
      % Fundo superior
      \fill[javaOrange]
        (current page.north west) rectangle
        ([yshift=-10cm] current page.north east);
      % Faixa decorativa
      \fill[javaDeep]
        ([yshift=-10cm] current page.north west) rectangle
        ([yshift=-10.5cm] current page.north east);
      % Faixa inferior
      \fill[javaAmber!40!white]
        (current page.south west) rectangle
        ([yshift=3cm] current page.south east);
    \end{tikzpicture}

    \vspace*{2.5cm}
    {\color{white}\fontsize{13}{16}\selectfont\textbf{PROGRAMAÇÃO EM JAVA}}\\[0.5cm]
    {\color{white}\fontsize{28}{34}\selectfont\textbf{Lição 3}}\\[0.3cm]
    {\color{white!90!javaOrange}\fontsize{18}{22}\selectfont Listas, Arrays e Strings}\\[1.8cm]

    {\color{white!85!javaOrange}\large Java Programming MOOC}\\[0.3cm]
    {\color{white!80!javaOrange}\normalsize Universidade de Helsinque}\\[2cm]

    \vfill
    {\color{javaAmber!60!black}\small
      Tradução para o português do Brasil — 2026\\
      Licença: CC BY-NC-SA 4.0
    }
    \vspace{1.5cm}
  \end{titlepage}
}

% ══════════════════════════════════════════════════════════
\begin{document}

\makejavacapa
\tableofcontents
\newpage

% ──────────────────────────────────────────────────────────
\section{Objetivos de aprendizagem}
% ──────────────────────────────────────────────────────────

Ao final desta lição, você será capaz de:

\begin{itemize}
  \item Usar \code{ArrayList} para armazenar coleções de dados
  \item Adicionar, acessar, remover e buscar elementos em uma lista
  \item Percorrer listas com laços \code{for} e \code{for-each}
  \item Criar e usar arrays de tamanho fixo
  \item Trabalhar com métodos importantes da classe \code{String}
  \item Dividir strings com \code{split()}
\end{itemize}

% ──────────────────────────────────────────────────────────
\section{Descobrindo e corrigindo erros}
% ──────────────────────────────────────────────────────────

Antes de falar sobre listas, é importante saber lidar com erros. Em Java, o compilador identifica erros de \textbf{sintaxe} (código mal escrito), mas erros de \textbf{lógica} só aparecem durante a execução.

\subsection{Tipos de erro}

\begin{center}
\rowcolors{2}{javaLightBg}{white}
\begin{longtable}{>{\bfseries}p{3.2cm} p{4.2cm} p{5.5cm}}
\toprule
\rowcolor{javaOrange}
\color{white}\textbf{Tipo} &
\color{white}\textbf{Quando ocorre} &
\color{white}\textbf{Exemplo} \\
\midrule
\endhead
Erro de sintaxe & Na compilação & Falta de \code{;} ou \code{\}} \\
Erro de execução & Ao rodar o programa & Divisão por zero, \code{NullPointerException} \\
Erro de lógica & O programa roda, mas dá resultado errado & Usar \code{+} em vez de \code{-} \\
\bottomrule
\end{longtable}
\end{center}

\subsection{Stack trace}

Quando um erro de execução ocorre, Java exibe um \textbf{stack trace} — uma lista das chamadas de método até o ponto do erro. Leia de baixo para cima para entender o que causou o problema:

\begin{plaincode}
Exception in thread "main" java.lang.ArrayIndexOutOfBoundsException: Index 5 out of bounds for length 5
    at Programa.main(Programa.java:8)
\end{plaincode}

Isso diz: o erro foi \code{ArrayIndexOutOfBoundsException} na linha 8 do arquivo \code{Programa.java}.

% ──────────────────────────────────────────────────────────
\section{Listas com \texttt{ArrayList}}
% ──────────────────────────────────────────────────────────

Um \code{ArrayList} é uma coleção dinâmica que pode crescer ou diminuir conforme necessário. É muito mais prático do que declarar dezenas de variáveis separadas.

\subsection{Criando e adicionando elementos}

\begin{javacode}
import java.util.ArrayList;

ArrayList<String> nomes = new ArrayList<>();
nomes.add("Alice");
nomes.add("Bruno");
nomes.add("Carlos");

System.out.println(nomes); // [Alice, Bruno, Carlos]
\end{javacode}

\begin{callout}[Dica]
Para tipos primitivos como \code{int} e \code{double}, use os equivalentes em maiúsculo:
\code{Integer}, \code{Double}.
\end{callout}

\begin{javacode}
ArrayList<Integer> numeros = new ArrayList<>();
numeros.add(10);
numeros.add(20);
numeros.add(30);
\end{javacode}

\subsection{Acessando elementos}

Os índices começam em 0:

\begin{javacode}
System.out.println(nomes.get(0)); // Alice
System.out.println(nomes.get(1)); // Bruno
System.out.println(nomes.size()); // 3
\end{javacode}

\subsection{Removendo elementos}

\begin{javacode}
nomes.remove("Bruno");        // remove pelo valor
System.out.println(nomes);    // [Alice, Carlos]

nomes.remove(0);              // remove pelo indice
System.out.println(nomes);    // [Carlos]
\end{javacode}

\subsection{Verificando existência}

\begin{javacode}
ArrayList<String> frutas = new ArrayList<>();
frutas.add("maca");
frutas.add("banana");

System.out.println(frutas.contains("banana")); // true
System.out.println(frutas.contains("uva"));    // false
\end{javacode}

\subsection{Percorrendo listas}

\subsubsection{Com \texttt{for} indexado}

\begin{javacode}
for (int i = 0; i < nomes.size(); i++) {
    System.out.println(i + ": " + nomes.get(i));
}
\end{javacode}

\subsubsection{Com \texttt{for-each} (mais legível)}

\begin{javacode}
for (String nome : nomes) {
    System.out.println(nome);
}
\end{javacode}

\subsubsection{Com \texttt{while}}

\begin{javacode}
int i = 0;
while (i < nomes.size()) {
    System.out.println(nomes.get(i));
    i++;
}
\end{javacode}

\subsection{Lista como parâmetro de método}

Como \code{ArrayList} é um tipo de referência, o método recebe a referência ao objeto original --- alterações dentro do método afetam a lista original:

\begin{javacode}
public static void adicionarZero(ArrayList<Integer> lista) {
    lista.add(0);
}

public static void main(String[] args) {
    ArrayList<Integer> numeros = new ArrayList<>();
    numeros.add(1);
    numeros.add(2);

    adicionarZero(numeros);
    System.out.println(numeros); // [1, 2, 0]
}
\end{javacode}

\subsection{Exemplo: soma e média de uma lista}

\begin{javacode}
public static double calcularMedia(ArrayList<Integer> lista) {
    if (lista.size() == 0) {
        return 0.0;
    }
    int soma = 0;
    for (int n : lista) {
        soma += n;
    }
    return (double) soma / lista.size();
}
\end{javacode}

% ──────────────────────────────────────────────────────────
\section{Arrays}
% ──────────────────────────────────────────────────────────

Um \textbf{array} é uma coleção de \textbf{tamanho fixo} do mesmo tipo. Diferentemente do \code{ArrayList}, o tamanho não pode ser alterado após a criação.

\subsection{Criando arrays}

\begin{javacode}
int[] numeros = new int[5];         // array de 5 inteiros, todos 0
String[] nomes = new String[3];     // array de 3 strings, todas null

// Inicializacao direta
int[] valores = {10, 20, 30, 40, 50};
\end{javacode}

\subsection{Acessando e modificando elementos}

\begin{javacode}
int[] arr = {5, 10, 15, 20};

System.out.println(arr[0]); // 5
System.out.println(arr[3]); // 20

arr[1] = 99;
System.out.println(arr[1]); // 99
\end{javacode}

\subsection{Tamanho do array}

\begin{javacode}
int[] arr = {1, 2, 3, 4, 5};
System.out.println(arr.length); // 5
\end{javacode}

\subsection{Percorrendo um array}

\begin{javacode}
for (int i = 0; i < arr.length; i++) {
    System.out.println(arr[i]);
}

// Ou com for-each:
for (int n : arr) {
    System.out.println(n);
}
\end{javacode}

\subsection{Quando usar array vs.\ ArrayList?}

\begin{center}
\rowcolors{2}{javaLightBg}{white}
\begin{longtable}{p{5cm} >{\centering\arraybackslash}p{2.5cm} >{\centering\arraybackslash}p{2.5cm}}
\toprule
\rowcolor{javaOrange}
\color{white}\textbf{Critério} &
\color{white}\textbf{Array} &
\color{white}\textbf{ArrayList} \\
\midrule
\endhead
Tamanho fixo & Sim & Não \\
Mais eficiente & Sim & Não \\
Mais flexível & Não & Sim \\
Tipos primitivos diretos & Sim & Não (usa wrapper) \\
\bottomrule
\end{longtable}
\end{center}

Use arrays quando o tamanho for fixo e conhecido. Prefira \code{ArrayList} para a maioria dos casos.

% ──────────────────────────────────────────────────────────
\section{Trabalhando com Strings}
% ──────────────────────────────────────────────────────────

\code{String} é uma das classes mais usadas em Java. Ela possui muitos métodos úteis.

\subsection{Métodos principais}

\begin{javacode}
String texto = "Ola, mundo!";

System.out.println(texto.length());          // 11 -- tamanho
System.out.println(texto.toUpperCase());     // OLA, MUNDO!
System.out.println(texto.toLowerCase());     // ola, mundo!
System.out.println(texto.contains("mundo")); // true
System.out.println(texto.startsWith("Ola")); // true
System.out.println(texto.endsWith("!"));     // true
System.out.println(texto.indexOf("mundo"));  // 5 -- posicao
System.out.println(texto.replace("mundo", "Java")); // Ola, Java!
System.out.println(texto.trim());            // remove espacos das bordas
\end{javacode}

\subsection{Acessando caracteres}

\begin{javacode}
String palavra = "Java";
char c = palavra.charAt(0); // J
System.out.println(c);
\end{javacode}

\subsection{Comparando strings}

\begin{callout}[Atenção]
Sempre use \code{.equals()}, \textbf{nunca} \code{==} para comparar strings!
\end{callout}

\begin{javacode}
String a = "Java";
String b = "Java";

System.out.println(a.equals(b));                // true
System.out.println(a.equalsIgnoreCase("java")); // true -- ignora maiusculas/minusculas
\end{javacode}

\subsection{Dividindo strings com \texttt{split()}}

O método \code{split()} divide uma string em partes, retornando um array de Strings:

\begin{javacode}
String linha = "Alice,25,Sao Paulo";
String[] partes = linha.split(",");

System.out.println(partes[0]); // Alice
System.out.println(partes[1]); // 25
System.out.println(partes[2]); // Sao Paulo
\end{javacode}

\begin{javacode}
String frase = "Java e uma linguagem poderosa";
String[] palavras = frase.split(" ");

for (String palavra : palavras) {
    System.out.println(palavra);
}
\end{javacode}

\subsection{Convertendo tipos}

\begin{javacode}
int numero = Integer.valueOf("42");
double decimal = Double.valueOf("3.14");
String texto = String.valueOf(100); // "100"
\end{javacode}

\subsection{Exemplo: processando dados CSV}

\begin{javacode}
import java.util.Scanner;
import java.util.ArrayList;

public class ProcessadorCSV {
    public static void main(String[] args) {
        Scanner scanner = new Scanner(System.in);

        System.out.println("Digite registros no formato 'nome,idade' (ou 'fim' para encerrar):");

        ArrayList<String> nomes = new ArrayList<>();
        ArrayList<Integer> idades = new ArrayList<>();

        while (true) {
            String linha = scanner.nextLine();
            if (linha.equals("fim")) {
                break;
            }
            String[] partes = linha.split(",");
            nomes.add(partes[0]);
            idades.add(Integer.valueOf(partes[1]));
        }

        System.out.println("\n--- Registros ---");
        for (int i = 0; i < nomes.size(); i++) {
            System.out.println(nomes.get(i) + " tem " + idades.get(i) + " anos.");
        }
    }
}
\end{javacode}

% ──────────────────────────────────────────────────────────
\section{Resumo}
% ──────────────────────────────────────────────────────────

\begin{itemize}
  \item \code{ArrayList} é uma coleção dinâmica; use \code{add()}, \code{get()}, \code{remove()}, \code{contains()}, \code{size()}.
  \item Índices começam em 0; acessar índice inválido lança \code{IndexOutOfBoundsException}.
  \item Arrays têm tamanho fixo, declarado com \code{new tipo[tamanho]} ou literal \code{\{v1, v2, ...\}}.
  \item \code{String} tem métodos como \code{length()}, \code{contains()}, \code{split()}, \code{equals()}, \code{charAt()}.
  \item Para comparar strings, use \code{.equals()} --- nunca \code{==}.
  \item \code{split(",")} divide uma string em partes usando o delimitador fornecido.
\end{itemize}

% ──────────────────────────────────────────────────────────
\section{Exercícios}
% ──────────────────────────────────────────────────────────

\begin{exercicio}{1 --- Lista de compras}
Escreva um programa que:
\begin{itemize}
  \item Lê itens de uma lista de compras até o usuário digitar \code{"sair"}
  \item Exibe todos os itens numerados
  \item Informa a quantidade total de itens
\end{itemize}
\end{exercicio}

\begin{exercicio}{2 --- Mínimo e máximo}
Crie um método \code{minimo(ArrayList<Integer> lista)} e \code{maximo(ArrayList<Integer> lista)}.
Popule uma lista com 6 números lidos do usuário e exiba o menor e o maior.
\end{exercicio}

\begin{exercicio}{3 --- Palavras únicas}
Leia palavras do usuário (até \code{"fim"}) e armazene apenas as que ainda não foram adicionadas à lista (sem duplicatas). Ao final, exiba todas as palavras únicas.
\end{exercicio}

\begin{exercicio}{4 --- Invertendo um array}
Crie um método \code{inverter(int[] arr)} que retorna um novo array com os elementos na ordem inversa. Não modifique o array original.
\end{exercicio}

\begin{exercicio}{5 --- Análise de frase}
Escreva um programa que lê uma frase e exibe:
\begin{itemize}
  \item Número de palavras
  \item A palavra mais longa
  \item A palavra mais curta
  \item A frase com todas as palavras em maiúsculo
\end{itemize}
\end{exercicio}

% ──────────────────────────────────────────────────────────
\vfill
\begin{center}
  \small\color{gray}
  Java Programming MOOC · Universidade de Helsinque ·
  \href{https://java-programming.mooc.fi/part-3}{java-programming.mooc.fi/part-3}\\
  Licença: \href{https://creativecommons.org/licenses/by-nc-sa/4.0/}{CC BY-NC-SA 4.0} ·
  Tradução para o português do Brasil --- 2026
\end{center}

\end{document}
